% GENERATED FILE. Edit canonical lessons, then run npm run generate:books.
% canonical-source-hash: fnv1a64:4c2197103cb56adc
% canonical-lessons: SA-C284-yadi-tarhi, SA-C284-yadi-varsa-bhavisyati, SA-C284-anyatha, SA-R284-if-and-otherwise

\chapter{If It Rains}
\label{ch:sa-if-it-rains}
\hlchaptermodality{284}

\begin{chapteropening}
\textbf{By the end of this chapter:} \emph{I can say what follows if something happens (yadi … tarhi …), with the present or the future in both halves, and what will happen otherwise (anyathā).}

Complete SA-R284-if-and-otherwise: say what follows if something happens, and what happens otherwise.
\end{chapteropening}

\section[\texorpdfstring{\sk{यदि} … \sk{तर्हि} … (yadi … tarhi …)}{यदि … तर्हि … (yadi … tarhi …)}]{\sk{यदि} … \sk{तर्हि} … (yadi … tarhi …) — if, then}

\label{lesson:SA-C284-yadi-tarhi}



\begin{quote}
Before the new one: say \textquotedblleft{}if\textquotedblright{} and \textquotedblleft{}then, in that case\textquotedblright{}. (\emph{yadi}, \emph{tarhi}.) You have both words; this chapter puts them to work.
\end{quote}

\subsection*{You'll want to know: \sk{यदि} … \sk{तर्हि}}
\textbf{\sk{यदि}} (\emph{yadi}), \textquotedblleft{}if\textquotedblright{}, opens the condition; \textbf{\sk{तर्हि}} (\emph{tarhi}), \textquotedblleft{}then\textquotedblright{}, opens what follows from it:

\begin{quote}
\textbf{\sk{यदि} \sk{वर्षा} \sk{भवति},}
\end{quote}

\begin{quote}
\emph{yadi varṣā bhavati,}
\end{quote}

\begin{quote}
\textbf{\sk{तर्हि} \sk{अहं} \sk{गृहे} \sk{तिष्ठामि।}}
\end{quote}

\begin{quote}
\emph{tarhi ahaṃ gṛhe tiṣṭhāmi.} — \textquotedblleft{}If it rains, I stay at home.\textquotedblright{}
\end{quote}

\begin{quote}
\textbf{\sk{यदि} \sk{सः} \sk{आगच्छति},}
\end{quote}

\begin{quote}
\emph{yadi saḥ āgacchati,}
\end{quote}

\begin{quote}
\textbf{\sk{तर्हि} \sk{वयं} \sk{खादामः।}}
\end{quote}

\begin{quote}
\emph{tarhi vayaṃ khādāmaḥ.} — \textquotedblleft{}If he comes, we eat.\textquotedblright{}
\end{quote}

\subsection*{You'll want to know: if not, and the condition second}
\textquotedblleft{}If not\textquotedblright{} puts \textbf{\sk{न}} before the verb of the \textbf{\sk{यदि}} half:

\begin{quote}
\textbf{\sk{यदि} \sk{सः} \sk{न} \sk{आगच्छति},}
\end{quote}

\begin{quote}
\emph{yadi saḥ na āgacchati,}
\end{quote}

\begin{quote}
\textbf{\sk{तर्हि} \sk{अहं} \sk{गच्छामि।}}
\end{quote}

\begin{quote}
\emph{tarhi ahaṃ gacchāmi.} — \textquotedblleft{}If he does not come, I go.\textquotedblright{}
\end{quote}

The condition can also come second, with no \textbf{\sk{तर्हि}}:

\begin{quote}
\textbf{\sk{अहं} \sk{गच्छामि}, \sk{यदि} \sk{सः} \sk{न} \sk{आगच्छति।}}
\end{quote}

\begin{quote}
\emph{ahaṃ gacchāmi, yadi saḥ na āgacchati.} — \textquotedblleft{}I go if he does not come.\textquotedblright{}
\end{quote}

\subsection*{Your turn}
Say these aloud:

\begin{itemize}
\raggedright
  \item \textquotedblleft{}If he comes, we eat\textquotedblright{} — \emph{yadi saḥ āgacchati, tarhi vayaṃ khādāmaḥ}
  \item \textquotedblleft{}If he comes, we read\textquotedblright{} — \emph{yadi saḥ āgacchati, tarhi vayaṃ paṭhāmaḥ}
  \item \textquotedblleft{}If he comes, we play\textquotedblright{} — \emph{yadi saḥ āgacchati, tarhi vayaṃ krīḍāmaḥ}
\end{itemize}

\begin{tcolorbox}[breakable,colback=teal!4,colframe=teal!35!black,title={Before you move on}]
\textquotedblleft{}If\textquotedblright{}? (\textbf{\sk{यदि}}.) \textquotedblleft{}Then\textquotedblright{}? (\textbf{\sk{तर्हि}}.) \textquotedblleft{}If he does not come\textquotedblright{}? (\textbf{\sk{यदि} \sk{सः} \sk{न} \sk{आगच्छति}}.)
\end{tcolorbox}
\section[\texorpdfstring{\sk{यदि} \sk{वर्षा} \sk{भविष्यति} (yadi varṣā bhaviṣyati)}{यदि वर्षा भविष्यति (yadi varṣā bhaviṣyati)}]{\sk{यदि} \sk{वर्षा} \sk{भविष्यति} (yadi varṣā bhaviṣyati) — if it happens, it will follow}

\label{lesson:SA-C284-yadi-varsa-bhavisyati}



\begin{quote}
Say \textquotedblleft{}if it rains, I stay at home\textquotedblright{}. (\emph{yadi varṣā bhavati, tarhi ahaṃ gṛhe tiṣṭhāmi}.) That is what happens every time. Now say what will happen tomorrow.
\end{quote}

\subsection*{You'll want to know: the future in both halves}
For one real thing ahead, both halves take the future, as Lanman's vocabulary shows for \textbf{\sk{यदि}}:

\begin{quote}
\textbf{\sk{यदि} \sk{वर्षा} \sk{भविष्यति},}
\end{quote}

\begin{quote}
\emph{yadi varṣā bhaviṣyati,}
\end{quote}

\begin{quote}
\textbf{\sk{तर्हि} \sk{अहं} \sk{गृहे} \sk{स्थास्यामि।}}
\end{quote}

\begin{quote}
\emph{tarhi ahaṃ gṛhe sthāsyāmi.} — \textquotedblleft{}If it rains, I will stay at home.\textquotedblright{}
\end{quote}

\begin{quote}
\textbf{\sk{यदि} \sk{समयः} \sk{भविष्यति},}
\end{quote}

\begin{quote}
\emph{yadi samayaḥ bhaviṣyati,}
\end{quote}

\begin{quote}
\textbf{\sk{तर्हि} \sk{अहम्} \sk{आगमिष्यामि।}}
\end{quote}

\begin{quote}
\emph{tarhi aham āgamiṣyāmi.} — \textquotedblleft{}If there is time, I will come.\textquotedblright{}
\end{quote}

\begin{quote}
\textbf{\sk{यदि} \sk{भवान्} \sk{आगमिष्यति},}
\end{quote}

\begin{quote}
\emph{yadi bhavān āgamiṣyati,}
\end{quote}

\begin{quote}
\textbf{\sk{तर्हि} \sk{मीरा} \sk{प्रसन्ना} \sk{भविष्यति।}}
\end{quote}

\begin{quote}
\emph{tarhi mīrā prasannā bhaviṣyati.} — \textquotedblleft{}If you come, Mira will be pleased.\textquotedblright{}
\end{quote}

\subsection*{You'll want to know: if …, then do this}
After a condition in the present, the second half can ask or advise, with \textbf{-\sk{तु}} or \textbf{-\sk{एत्}}:

\begin{quote}
\textbf{\sk{यदि} \sk{भवान्} \sk{रुग्णः} \sk{अस्ति},}
\end{quote}

\begin{quote}
\emph{yadi bhavān rugṇaḥ asti,}
\end{quote}

\begin{quote}
\textbf{\sk{तर्हि} \sk{वैद्यं} \sk{गच्छतु।}}
\end{quote}

\begin{quote}
\emph{tarhi vaidyaṃ gacchatu.} — \textquotedblleft{}If you are ill, please go to the doctor.\textquotedblright{}
\end{quote}

\begin{quote}
\textbf{\sk{यदि} \sk{धूमयानं} \sk{गतम्},}
\end{quote}

\begin{quote}
\emph{yadi dhūmayānaṃ gatam,}
\end{quote}

\begin{quote}
\textbf{\sk{तर्हि} \sk{भवान्} \sk{लोकयानेन} \sk{गच्छेत्।}}
\end{quote}

\begin{quote}
\emph{tarhi bhavān lokayānena gacchet.} — \textquotedblleft{}If the train has gone, you should go by bus.\textquotedblright{}
\end{quote}

\textbf{\sk{लोकयानेन}} (\emph{lokayānena}) is \textquotedblleft{}by bus\textquotedblright{}: \emph{-ena}, \textquotedblleft{}by\textquotedblright{}, as in \textbf{\sk{मतेन}}.

\subsection*{Your turn}
Say these aloud:

\begin{itemize}
\raggedright
  \item \textquotedblleft{}If there is time, I will come\textquotedblright{} — \emph{yadi samayaḥ bhaviṣyati, tarhi aham āgamiṣyāmi}
  \item \textquotedblleft{}If there is time, I will read\textquotedblright{} — \emph{yadi samayaḥ bhaviṣyati, tarhi ahaṃ paṭhiṣyāmi}
  \item \textquotedblleft{}If there is time, I will write\textquotedblright{} — \emph{yadi samayaḥ bhaviṣyati, tarhi ahaṃ lekhiṣyāmi}
\end{itemize}

\begin{tcolorbox}[breakable,colback=teal!4,colframe=teal!35!black,title={Before you move on}]
\textquotedblleft{}If it rains tomorrow, I will stay at home\textquotedblright{}? (\textbf{\sk{यदि} \sk{श्वः} \sk{वर्षा} \sk{भविष्यति}, \sk{तर्हि} \sk{अहं} \sk{गृहे} \sk{स्थास्यामि}}.) Which form goes in both halves? (The future.) \textquotedblleft{}If you are ill, please go to the doctor\textquotedblright{}? (\textbf{\sk{यदि} \sk{भवान्} \sk{रुग्णः} \sk{अस्ति}, \sk{तर्हि} \sk{वैद्यं} \sk{गच्छतु}}.)
\end{tcolorbox}
\section[\texorpdfstring{\sk{अन्यथा} (anyathā)}{अन्यथा (anyathā)}]{\sk{अन्यथा} (anyathā) — otherwise}

\label{lesson:SA-C284-anyatha}



\begin{quote}
Before the new one: say \textquotedblleft{}if there is time, I will come\textquotedblright{}. (\emph{yadi samayaḥ bhaviṣyati, tarhi aham āgamiṣyāmi}.) Then say \textquotedblleft{}quickly\textquotedblright{}. (\emph{śīghram}.)
\end{quote}

\subsection*{You'll want to know: \sk{अन्यथा}}
\textbf{\sk{अन्यथा}} — \emph{anyathā} — \textquotedblleft{}otherwise\textquotedblright{}: if not, then this.

\begin{quote}
\textbf{\sk{शीघ्रं} \sk{गच्छतु},}
\end{quote}

\begin{quote}
\emph{śīghraṃ gacchatu,}
\end{quote}

\begin{quote}
\textbf{\sk{अन्यथा} \sk{धूमयानं} \sk{गमिष्यति।}}
\end{quote}

\begin{quote}
\emph{anyathā dhūmayānaṃ gamiṣyati.} — \textquotedblleft{}Go quickly, otherwise the train will leave.\textquotedblright{}
\end{quote}

\begin{quote}
\textbf{\sk{कृपया} \sk{शीघ्रम्} \sk{आगच्छतु},}
\end{quote}

\begin{quote}
\emph{kṛpayā śīghram āgacchatu,}
\end{quote}

\begin{quote}
\textbf{\sk{अन्यथा} \sk{अहं} \sk{गमिष्यामि।}}
\end{quote}

\begin{quote}
\emph{anyathā ahaṃ gamiṣyāmi.} — \textquotedblleft{}Please come quickly, otherwise I will leave.\textquotedblright{}
\end{quote}

\begin{cousinweb}
\textbf{\sk{अन्यथा}} (\emph{anyathā}) is built on \textbf{\sk{अन्य}} (\emph{anya}), \textquotedblleft{}other\textquotedblright{}, with \emph{-thā}, \textquotedblleft{}in the way of\textquotedblright{}: \textquotedblleft{}in another way\textquotedblright{}. Monier-Williams gives it as \textquotedblleft{}otherwise, in a different manner\textquotedblright{}, and also \textquotedblleft{}in the contrary case\textquotedblright{}, the sense here.
\end{cousinweb}

\subsection*{Your turn}
Say these aloud:

\begin{itemize}
\raggedright
  \item \emph{anyathā}
  \item \textquotedblleft{}Go quickly, otherwise the train will leave\textquotedblright{} — \emph{śīghraṃ gacchatu, anyathā dhūmayānaṃ gamiṣyati}
  \item \textquotedblleft{}You should take the medicine, otherwise you will be ill\textquotedblright{} — \emph{bhavān auṣadhaṃ pibet, anyathā rugṇaḥ bhaviṣyati}
\end{itemize}

\begin{tcolorbox}[breakable,colback=teal!4,colframe=teal!35!black,title={Before you move on}]
What does \textbf{\sk{अन्यथा}} mean? (\textquotedblleft{}Otherwise\textquotedblright{}.) \textquotedblleft{}Otherwise I will leave\textquotedblright{}? (\textbf{\sk{अन्यथा} \sk{अहं} \sk{गमिष्यामि}}.)
\end{tcolorbox}
\section[(review)]{Review — if, and otherwise}

\label{lesson:SA-R284-if-and-otherwise}



\begin{quote}
Say \textquotedblleft{}the train goes faster than the bus\textquotedblright{}. (\emph{dhūmayānaṃ lokayānāpekṣayā śīghraṃ gacchati}.) Then \textquotedblleft{}from morning until evening\textquotedblright{}. (\emph{prātaḥ ārabhya sāyaṃ paryantam}.)
\end{quote}

\subsection*{Your turn}
Say these aloud:

\begin{itemize}
\raggedright
  \item If it rains, I stay at home.
  \item If he does not come, we will go.
  \item If there is time tomorrow, I will see the city.
  \item If you are tired, please sit down.
  \item Come quickly, otherwise the bus will leave.
\end{itemize}

\begin{tcolorbox}[breakable,colback=teal!4,colframe=teal!35!black,title={Before you move on}]
\textbf{\sk{यदि}} or \textbf{\sk{तर्हि}}, to open the condition? (\textbf{\sk{यदि}}.) For tomorrow, which form in both halves? (The future.) Now the other way round: what does \textbf{\sk{अन्यथा}} add? (\textquotedblleft{}Otherwise\textquotedblright{}: what happens if not.)
\end{tcolorbox}
